% Generated draft: requires mathematical and source review.
\documentclass[11pt]{article}
\usepackage[margin=25mm]{geometry}
\usepackage{fontspec}
\setmainfont{Latin Modern Roman}
\usepackage{amsmath,amssymb,mathtools,unicode-math}
\setmathfont{Latin Modern Math}
\usepackage{enumitem,needspace,xurl,hyperref}
\setlength{\parindent}{0pt}
\setlength{\parskip}{5pt}
\setlength{\emergencystretch}{3em}
\newcommand{\N}{\mathbb{N}}
\newcommand{\Z}{\mathbb{Z}}
\newcommand{\Q}{\mathbb{Q}}
\newcommand{\R}{\mathbb{R}}
\newcommand{\C}{\mathbb{C}}
\newcommand{\F}{\mathbb{F}}
\newcommand{\A}{\mathbb{A}}
\newcommand{\PP}{\mathbb{P}}
\newcommand{\E}{\mathbb{E}}
\newcommand{\eps}{\varepsilon}
\newcommand{\dd}{\,\mathrm{d}}
\begin{document}
\section*{UnsolvedMath ID: 801015}
\textbf{Source problem code:} catalogue entry
\subsection{Definitions and mathematical statement}
The supplied record does not define \(\mathbf{WF}\), \(\mathbf{UB}\), their ambient space, or the reduction notion used for completeness. Consequently, a fully specified mathematical statement cannot be recovered from this record alone.

Here is a conditional formulation under the usual boldface, continuous-reduction interpretation. Let \(X\) be the intended Polish space, and let \(W,U\subseteq X\) denote the intended sets \(\mathbf{WF}\) and \(\mathbf{UB}\). A subset of a Polish space is analytic if it is the projection of a closed subset of its product with Baire space \(\N^{\N}\); it is coanalytic if its complement is analytic. A coanalytic separator of \(W\) from \(U\) is a coanalytic set \(S\subseteq X\) satisfying
\[
W\subseteq S,\qquad S\cap U=\varnothing.
\]
Under continuous reductions from Baire space, a coanalytic set \(S\subseteq X\) is coanalytic complete if, for every coanalytic set \(B\subseteq\N^{\N}\), there exists a continuous map \(f:\N^{\N}\to X\) with \(B=f^{-1}(S)\).

\textbf{Conditional statement and requested proof method.} Give a proof using descriptive set theoretic methods that
\[
\forall S\subseteq X,\quad
\bigl[S\text{ is coanalytic}\ \land\ W\subseteq S\ \land\ S\cap U=\varnothing\bigr]
\Longrightarrow
\bigl[S\text{ is coanalytic complete}\bigr].
\]
The source attributes a recursion theoretic proof of the intended assertion to Goh and states that the same method also proves Saint Raymond's theorem. It asks whether the assertion can be proved using descriptive set theoretic methods. The particular Saint Raymond theorem is not stated, and whether the question also requests such a proof of that theorem is unclear. The displayed formulation remains conditional: the missing definitions and conventions must be recovered before it can serve as the intended precise statement.
\subsection{Short English statement}
The source asks for a descriptive set theoretic proof that every coanalytic separator of \(\mathbf{WF}\) and \(\mathbf{UB}\) is coanalytic complete, attributing a recursion theoretic proof to Goh. The sets and completeness convention are unspecified, so the mathematical formulation above is conditional. The supplied registry labels the problem open; that label has not been independently verified.
\subsection{Sources}
\begin{itemize}
\item Supplied local UnsolvedMath registry record: permanent ID \(801015\), exact source problem code \texttt{catalogue entry}, title ``Open Problems in Mathematical Logic --- Problem 27.'' The record lists Barmpalias, Gao et al. as proposers and gives the year 2026.
\item Supplied source attribution: Barmpalias, Gao et al., \emph{Open Problems in Mathematical Logic} (2026), Problem 27, \url{https://arxiv.org/abs/2608.26628}. The background records ``Accessed: 2026-10-16'' and extraction method \texttt{arxiv\_latex}; these are supplied provenance claims, not independently verified access.
\item The statement attributes the recursion theoretic result to Goh and mentions Saint Raymond's theorem, but supplies no separate bibliographic references or theorem statements for either attribution. The cited source list has not been independently consulted.
\end{itemize}
\end{document}
